% TOP_PROBLEM: {"id": 30001796, "title": "Algebraic K-Theory and Rigidity Conjectures for Discrete Groups", "category_id": 17, "category": {"id": 17, "name": "group_theory", "display_name": "Group Theory", "description": "Problems about groups, group actions, representations, and related algebraic structures.", "slug": "group-theory", "order_index": 17, "created_at": "2026-07-31T15:26:25.671Z"}, "status": "partially_solved"}
% FIRST_POSTED: null
% ENTRY_KIND: statement_only
% !TeX program = lualatex
\documentclass[11pt]{article}
\usepackage[margin=25mm]{geometry}
\usepackage{fontspec}
\setmainfont{Latin Modern Roman}
\usepackage{amsmath,amssymb,mathtools}
\usepackage{unicode-math}
\setmathfont{Latin Modern Math}
\usepackage{xurl,hyperref}
\hypersetup{colorlinks=true,urlcolor=blue}
\setcounter{secnumdepth}{0}
\setlength{\parindent}{0pt}
\setlength{\parskip}{5pt}
\setlength{\emergencystretch}{3em}
\newcommand{\R}{\mathbb R}
\newcommand{\N}{\mathbb N}
\newcommand{\Z}{\mathbb Z}
\newcommand{\Q}{\mathbb Q}
\newcommand{\C}{\mathbb C}
\newcommand{\D}{\mathbb D}
\newcommand{\F}{\mathbb F}
\newcommand{\sref}[2]{\hyperlink{source:#1:#2}{[S#2]}}
\newcommand{\cataloguescope}[1]{#1}
\begin{document}
% UnsolvedMath ID 30001796; source catalogue code OWR-5155-002
\setcounter{section}{30001795}
\section{Algebraic K-Theory and Rigidity Conjectures for Discrete Groups}\label{30001796}
{\small UnsolvedMath ID 30001796 · Group Theory}
\subsection{Definitions and mathematical statement}

\paragraph{Shared notation.}$\mathbb N$, $\mathbb Z$, $\mathbb Q$, $\mathbb R$, and $\mathbb C$ denote the natural numbers, integers, rationals, reals, and complex numbers, respectively. All additional parameters and hypotheses are specified in the question.


\paragraph{Statement recorded in the supplied source.}
For a torsion-free group $G$ and an integral domain $R$, are $0$ and $1$ the only idempotents of $RG$? Do the negative K-groups of $\mathbb{Z}G$, the reduced group $K_0(\mathbb{Z}G)$, and $\operatorname{Wh}(G)$ vanish? Are higher signatures homotopy invariant, closed aspherical manifolds topologically rigid, and groups of type FP of type FF? Is every finitely presented Poincaré-duality group of dimension at least five the fundamental group of an aspherical homology ANR-manifold? Is every hyperbolic group with boundary $S^n$ the fundamental group of a closed aspherical manifold? Are the Farrell–Jones assembly maps in algebraic K- and L-theory bijective?
\par\smallskip\textit{Formulation references:} \sref{30001796}{1}, \sref{30001796}{2}.
\subsection{Short English statement}
For a torsion-free group $G$ and an integral domain $R$, are $0$ and $1$ the only idempotents of $RG$? Do the negative K-groups of $\mathbb{Z}G$, the reduced group $K_0(\mathbb{Z}G)$, and $\operatorname{Wh}(G)$ vanish? Are higher signatures homotopy invariant, closed aspherical manifolds topologically rigid, and groups of type FP of type FF? Is every finitely presented Poincaré-duality group of dimension at least five the fundamental group of an aspherical homology ANR-manifold? Is every hyperbolic group with boundary $S^n$ the fundamental group of a closed aspherical manifold? Are the Farrell–Jones assembly maps in algebraic K- and L-theory bijective?
\subsection{Sources}
\begin{itemize}
\item[\hypertarget{source:30001796:1}{S1}] ULAM.AI and the UnsolvedMath Contributors, UnsolvedMath: A Curated Collection of Open Mathematics Problems, record 30001796 (OWR-5155-002). Catalogue attribution under CC BY 4.0.
\url{https://huggingface.co/datasets/ulamai/UnsolvedMath}.
\item[\hypertarget{source:30001796:2}{S2}] Original problem formulation and mathematical context.
\url{https://doi.org/10.4171/owr/2011/27}.
\end{itemize}
\label{end:30001796}
\end{document}
